\section{Persistent guarded segments and the same pointwise target}
\label{sec:v110-persistence}
The positive decomposition of the preceding section has a geometric
consequence on every segment which retains source. It uses the
original incidence and clearance guard, without adding a distance
to the section boundary.

\begin{theorem}[Persistence of the guard on the remaining source]
\label{thm:v110-persistence}
Let $\gamma>1$. If $c_{\rm per}^{\varepsilon,\gamma}(\theta)>0$,
then $c^\varepsilon>0$ and, at every point of the closed physical
segment,
\begin{equation}\label{eq:v110-near-constant}
 |f^\varepsilon(u,\theta)-c^\varepsilon(\theta)|
       \le\mathcal V^\varepsilon(\theta)
       <\gamma^{-1}c^\varepsilon(\theta).
\end{equation}
In particular $f^\varepsilon>0$ and $W^\varepsilon>0$ throughout
that segment, including its one-sided endpoints. For each $u$ at
least one of the original tests satisfies
\begin{equation}\label{eq:v110-witness}
 c_j(u)<2\varepsilon q_0^{d_j/4}
       \quad\text{for some }0\le j\le m,
 \qquad\text{or}\qquad
 \Delta_j(u)<2\varepsilon q_0^{d'_j/4}
       \quad\text{for some }0\le j<m.
\end{equation}
The witnessing test may change with $u$.

Conversely, if $W^\varepsilon$ vanishes anywhere on the closed
segment, its complete original paired coefficient is paid:
$c_{\rm var}^{\varepsilon,\gamma}=c^\varepsilon$ and
$c_{\rm per}^{\varepsilon,\gamma}=0$.
\end{theorem}
\begin{proof}
The endpoint minimum is attained at at least one endpoint, where
$f^\varepsilon=c^\varepsilon$. Every increment from that endpoint
is bounded by the total variation. Positivity of the remaining
coefficient is precisely $c^\varepsilon>\gamma\mathcal V^\varepsilon$,
which proves \eqref{eq:v110-near-constant}. If any visited interior
incidence vanished at an endpoint, Theorem~\ref{thm:v109-traces}
would give $c^\varepsilon=0$. Thus the segment belongs to
$\mathcal N$. Its optical factor $A$ is finite and positive on the
closed segment, including a critical endpoint with its fixed
Gauss direction. Since $f^\varepsilon=AW^\varepsilon>0$, the guard
is positive there. If all inequalities in \eqref{eq:v110-witness}
failed, every factor defining $H^\varepsilon$ would equal one,
contradicting $W^\varepsilon=1-H^\varepsilon>0$.

For the converse, let $u_0$ be a point where $W^\varepsilon=0$.
The finite optical factor gives $f^\varepsilon(u_0)=0$, so the
increment from an endpoint attaining $c^\varepsilon$ gives
$\mathcal V^\varepsilon\ge c^\varepsilon$. Apply
\eqref{eq:v110-positive-split}. This argument also covers $u_0=a$
or $u_0=b$, with the proved physical one-sided values.
\end{proof}

The surviving class
\begin{equation}\label{eq:v110-persistent-class}
 \mathcal P_B=\{\theta\in\mathcal N:
       c^{\varepsilon(B)}(\theta)>
                      \gamma_B\mathcal V^{\varepsilon(B)}(\theta)\}
\end{equation}
is Borel. The source $\mathsf U_B^{\rm per}$ is supported on it;
its coefficient is still the positive difference in
\eqref{eq:v110-positive-split}, not the full original coefficient.
The sum \eqref{eq:v110-source-identity} retains every original
coefficient exactly. No assumption that one fixed bad test covers
the entire segment is made. A Borel witness at a specified point
can be chosen as the first successful test in the finite ordered
list in \eqref{eq:v110-witness}; this selection is not needed to
prove the payment.

\begin{remark}[Boundary mechanisms and what is retained]
An unvisited third-disk tangency may give $W^\varepsilon=1$ at an
endpoint without forcing variation relative to the endpoint minimum.
It is therefore not declared paid merely on account of its name.
The same applies to a nongrazing critical entrance, an original
section cut and a short-reach segment. Each receives precisely the
two coefficients in \eqref{eq:v110-positive-split}. When the guard
vanishes elsewhere on that segment the entire coefficient is paid
by Theorem~\ref{thm:v110-persistence}; otherwise a positive remaining
coefficient is possible. At fixed collision count the optical data
on a surviving segment have positive minima, but no such minimum
is asserted uniformly in the collision count or the sheet.
\end{remark}

\subsection{Comparison with the original nongrazing shortage}
Keep the \emph{full original} nongrazing pool $P_{r_B}^{\rm ng}$,
including its capacity on segments whose source has already been
paid. Define only a new source-side remainder:
\begin{equation}\label{eq:v110-persistent-residual}
 \mathsf R_B^{\rm per}(t)
   =\bigl(\mathsf U_B^{\rm per}(t)
                         -K_BP_{r_B}^{\rm ng}(t)\bigr)_+.
\end{equation}
The allocation by this pool is scalar at the roof, just as in
Section~\ref{sec:v109-shortage}. It is not a coupling of individual
persistent sheets to individual comparison sheets.

\begin{theorem}[Positive-source reduction with a paid error]
\label{thm:v110-residual}
On the original exact event, at every finite collision count,
\begin{align}
 \mathsf U_B^{\rm micro}
  &=\mathsf U_B^{\rm var}
    +\min\{\mathsf U_B^{\rm per},K_BP_{r_B}^{\rm ng}\}
    +\mathsf R_B^{\rm per},\label{eq:v110-three-parts}\\
 0\le\mathsf R_B^{\rm ng}-\mathsf R_B^{\rm per}
  &\le\mathsf U_B^{\rm var}.\label{eq:v110-shortage-difference}
\end{align}
For every fixed $M$, in the prescribed order of limits,
\begin{align}
 \mathcal H_M(\mathsf R_B^{\rm ng}-\mathsf R_B^{\rm per})
     &\le C_MB^{-95/384},\label{eq:v110-height-difference}\\
 \mathcal E_M&\le\mathcal H_M(\mathsf R_B^{\rm per})
                           +C_MB^{-1/384}.
                           \label{eq:v110-final-reduction}
\end{align}
The unchanged pointwise assertion \eqref{eq:v76-target} is equivalent
to $\mathcal H_M(\mathsf R_B^{\rm per})\to0$ for every fixed $M$.
\end{theorem}
\begin{proof}
Equation \eqref{eq:v110-three-parts} is the positive scalar identity
$x=\min(x,y)+(x-y)_+$ applied to $x=\mathsf U_B^{\rm per}$,
combined with \eqref{eq:v110-source-identity}. Monotonicity and the
unit Lipschitz bound of the positive-part function give
\eqref{eq:v110-shortage-difference}. Corollary~\ref{cor:v110-schedule}
then proves \eqref{eq:v110-height-difference}.

For completeness, the height of the middle term in
\eqref{eq:v110-three-parts} is at most
\[
 \mathcal H_M(K_BP_{r_B})
 \le C\varepsilon^{-1/32}
       \{\varepsilon^{1/16}+(1+\varepsilon^{-1})\Delta_B\}
 =O(\varepsilon^{1/32})=O(B^{-1/384}),
\]
by Theorem~\ref{thm:v109-pool}. Insert the three nonnegative terms
in Input statement~\ref{input:v109-scalar}, and use subadditivity
of the ordered upper height. This proves
\eqref{eq:v110-final-reduction} and sufficiency of the displayed
limit. Necessity follows from the same inherited scalar equivalence
and $0\le\mathsf R_B^{\rm per}\le\mathsf U_B^{\rm micro}$.
No essential supremum is interchanged with an integral, and the
comparison pool has not been restricted to $\mathcal P_B$.
\end{proof}

The estimate still required to conclude \eqref{eq:v76-target} is
\begin{equation}\label{eq:v110-remaining-limit}
 \lim_{B\to\infty}\limsup_{m\to\infty}\sup_{R,n,k}m^2
 \esssup_{t\in\mathcal C_{m,R,n,k}(M)}
 \bigl(\mathsf U_{B,m}^{\rm per}(t)
                    -K_BP_{r_B,m}^{\rm ng}(t)\bigr)_+=0.
\end{equation}
Theorems~\ref{thm:v110-payment} and \ref{thm:v110-persistence}
prove an actual summed payment and a necessary physical persistence
property. They do not prove \eqref{eq:v110-remaining-limit}. In
particular a positive nearly constant density on a shrinking segment
need not have small pointwise height merely because its roof integral
is small. The original raw pointwise theorem, and not an averaged
substitute, remains the organizing target.
